% Pandoc LaTeX template — arXiv preprint (article + arxiv.sty)
\documentclass[11pt]{article}

\usepackage{arxiv}
\usepackage{fontspec}
\defaultfontfeatures{Ligatures=TeX}
\setmainfont{lmroman10-regular.otf}[
  BoldFont=lmroman10-bold.otf,
  ItalicFont=lmroman10-italic.otf,
  BoldItalicFont=lmroman10-bolditalic.otf
]
\setmonofont{lmroman10-regular.otf}

\usepackage{amsmath,amssymb}
\usepackage{booktabs}
\usepackage{longtable}
\usepackage{array}
\usepackage{calc}
\usepackage{microtype}
\usepackage{hyperref}
\usepackage{url}
\usepackage{xcolor}
\usepackage{graphicx}
\usepackage{caption}
\usepackage{ltcaption}
\usepackage{enumitem}
\usepackage{framed}

\newcounter{none}
\newcommand{\real}[1]{#1}

\definecolor{shadecolor}{RGB}{248,248,248}
\newenvironment{Shaded}{\begin{snugshade}}{\end{snugshade}}

\providecommand{\tightlist}{%
  \setlength{\itemsep}{2pt}%
  \setlength{\parskip}{0pt}%
}

\hypersetup{
  colorlinks=true,
  linkcolor=black,
  citecolor=black,
  urlcolor=black,
  pdftitle={Domain-Oriented Context Cores: Bounded Context for AI
Decision Support in Software Projects},
  pdfauthor={Alexander Young}
}

\title{Domain-Oriented Context Cores: Bounded Context for AI Decision
Support in Software Projects}
\author{Alexander Young}
\date{June 2026}

\begin{document}

\maketitle

\begin{abstract}
Modern AI-assisted software development often loads repository-scale
context into large language models (LLMs), increasing cost, latency, and
attention dilution on scoped decision questions. We introduce AI Context
OS, a framework that models repository knowledge as Context Cores -
domain-oriented context partitions with lightweight multi-core routing.
We report exploratory within-subjects A/B/C experiments on four
codebases (149 decision-scoped questions; gpt-4o-mini; LLM-as-judge):
MailAgent, Django REST Framework, Navorina, and Oiloop (private macOS).
Context Cores (Condition B) achieved 8× to 38× token compression
(CCR\_tokens) and large cost reductions versus full-repository baselines
(Condition A). On three open-source projects, B mean accuracy exceeded A
by 19-24 percentage points. On Oiloop (Phase 3.1, 20 questions, v1.1
cores), B reached 2.70 versus A 0.75 (bootstrap 95\% CI on paired Δ:
{[}+1.60, +2.30{]}); a production keyword router preserved gains (B
2.75, F1 1.0, hallucination 0\%). Hybrid core-plus-graph retrieval
(Condition D) did not beat multi-core B on cross-cutting questions
(H\(_{1f}\) rejected). The primary hypothesis is supported on
decision-scoped questions: efficiency gains are consistent; accuracy
gains require adequate core metadata and multi-core routing on
integrated codebases. A 10-question author-rated blind pilot preferred B
in 70\% of cases, with the remaining 30\% rated equal (n=10; not
independent experts). Masked decode preference on prior runs: 75\%
B-or-equal. We release protocols, raw runs (three of four full; Django
aggregate-only), and routing artifacts for replication.
\end{abstract}

\keywords{Context engineering \and Large language models \and Software
engineering \and Retrieval-augmented generation \and Domain-driven
design}

\section{Introduction}\label{introduction}

The integration of large language models (LLMs) into software
engineering workflows has progressed from function-level code completion
to repository-level understanding and reasoning. Contemporary commercial
LLMs expose context windows of up to two million tokens, encouraging a
direct engineering response: loading the entire workspace - source
files, dependencies, documentation, and configuration - into the prompt
baseline.

This brute-force approach simplifies retrieval pipelines but introduces
three systematic costs:

\begin{enumerate}
\def\labelenumi{\arabic{enumi}.}
\tightlist
\item
  \textbf{Token cost inflation.} Processing tens of thousands of tokens
  per query is financially unsustainable at high query volumes.
\item
  \textbf{Latency degradation.} Inference time scales with prompt
  length, introducing multi-second delays that disrupt developer focus.
\item
  \textbf{Attention dilution.} LLMs frequently fail to locate and
  prioritize critical details when surrounded by large volumes of
  irrelevant code, yielding incorrect recommendations, omissions, or
  hallucinations {[}1{]}.
\end{enumerate}

To address these limitations, we propose \textbf{Context Cores}, a
structured alternative grounded in domain-driven design (DDD) {[}3{]}.
Instead of treating a repository as a monolithic file index, we
construct minimal, self-contained semantic units aligned with distinct
architectural responsibilities. A lightweight query router selects
active cores for each query, compiling only essential guidelines and
code symbols.

We present the \textbf{AI Context OS} framework and report a multi-phase
evaluation across three open-source codebases and one proprietary macOS
application. We compare:

\begin{itemize}
\tightlist
\item
  \textbf{Condition A (full-repository baseline):} The entire repository
  structure and files loaded into the LLM context.
\item
  \textbf{Condition B (routed Context Cores):} Scoped domain guidelines
  and relevant entities loaded via keyword or semantic routing.
\item
  \textbf{Condition C (graph-based retrieval):} A dependency graph built
  via static AST analysis, loading nodes and code snippets within two
  hops of query seeds.
\end{itemize}

\subsection{Contributions}\label{contributions}

This paper makes the following contributions:

\begin{enumerate}
\def\labelenumi{\arabic{enumi}.}
\tightlist
\item
  \textbf{AI Context OS}, a domain-oriented context framework with
  explicit cores, routing, and prompt compilation.
\item
  A \textbf{within-subjects evaluation protocol} spanning 149
  decision-scoped questions on four codebases, with published raw runs
  and replication artifacts.
\item
  \textbf{Empirical evidence} that routed Context Cores compress input
  tokens by 8×-38× while matching or exceeding full-repository accuracy
  when core metadata and routing are adequate.
\item
  A \textbf{stratum-dependent production decision matrix} comparing
  Conditions B, C, and D under stated limitations (LLM-as-judge, partial
  Django export).
\end{enumerate}

\section{Problem Formulation}\label{problem-formulation}

\subsection{Attention Dilution in Large Context
Windows}\label{attention-dilution-in-large-context-windows}

Let \(P\) denote the user query and \(R = \{f_1, f_2, \dots, f_n\}\) the
set of repository files. Under Condition A, the input prompt is

\[I_A = [\text{SystemInstructions}, R, P].\]

As \(|R|\) grows, token length \(L(I_A)\) increases linearly. Prior work
shows that retrieval performance over long prompts follows a U-shaped
curve: information in the middle of the context is frequently
under-attended {[}1{]}.

\subsection{File-Centric versus Decision-Centric
Organization}\label{file-centric-versus-decision-centric-organization}

Traditional retrieval systems (lexical RAG, semantic chunking) operate
at file granularity {[}4{]}. Developer queries are often
decision-centric. Answering \emph{``How do we restrict API write
operations on viewsets to JWT-authenticated users?''} does not require
every serializer or view class; it requires rules governing
serialization, view configuration, and authentication, plus specific
file routes.

\subsection{Observed Failure Modes}\label{observed-failure-modes}

In audits of naive coding assistants, we observed recurring failures on:

\begin{itemize}
\tightlist
\item
  \textbf{System API boundaries:} Confusion in helper utilities while
  missing event-delegation patterns.
\item
  \textbf{Settings configuration maps:} Outdated properties due to
  absent single sources of truth.
\item
  \textbf{Cross-cutting rules:} Hallucinated integration paths when only
  one subsystem appears in context.
\end{itemize}

\section{Related Work}\label{related-work}

\subsection{Retrieval-Augmented
Generation}\label{retrieval-augmented-generation}

Lexical retrieval (e.g., BM25) and dense embedding retrieval are widely
used to prune LLM contexts {[}4{]}. Naive RAG often chunks files into
arbitrary blocks, disrupting control flow. On codebases, this fragments
class declarations from method implementations, forcing models to infer
missing interfaces.

\subsection{Knowledge Graphs and
GraphRAG}\label{knowledge-graphs-and-graphrag}

Graph-based methods model code as entities and relations (calls,
inheritance, imports) {[}5, 6{]}. Graphs preserve structure but optimize
connectivity over sufficiency. A developer modifying a settings
parameter may receive a large adjacent subgraph of consumers, flooding
the prompt with irrelevant code.

\subsection{Code Intelligence and AST
Indexing}\label{code-intelligence-and-ast-indexing}

Static analysis assistants parse abstract syntax trees (ASTs) for
precise declaration maps. AST indices lack business and operational
semantics: architectural constraints expressed in prose (e.g.,
\emph{``never send emails without user confirmation''}) are not
captured.

\subsection{Context Engineering and
Compression}\label{context-engineering-and-compression}

Prompt pruning, token compression, and context distillation remove
redundant tokens post hoc {[}7{]}. These methods do not restructure
repository knowledge; token and latency costs re-emerge as repositories
scale.

\section{Hypotheses}\label{hypotheses}

We formalize the framework assumptions as follows:

\begin{itemize}
\tightlist
\item
  \textbf{Primary hypothesis (\(H_1\)):} Dynamically routed,
  domain-oriented Context Cores produce responses equal or superior in
  correctness, completeness, and actionability to full-repository
  contexts, while significantly reducing input tokens and latency.
\item
  \textbf{\(H_{1a}\) (token reduction):} \(L(I_B) \ll L(I_A)\).
\item
  \textbf{\(H_{1b}\) (latency reduction):} \(t_B \ll t_A\).
\item
  \textbf{\(H_{1c}\) (domain accuracy):} On business-rule and
  domain-logic questions, \(\text{Acc}_B \ge \text{Acc}_A\).
\item
  \textbf{\(H_{1d}\) (hallucination mitigation):} Context Cores reduce
  structural hallucinations via explicit boundaries and exclusions.
\item
  \textbf{\(H_{1e}\) (expert preference):} Under blinded preference
  rating, domain experts prefer or rate as equal Condition B in at least
  60\% of cases versus Condition A.
\end{itemize}

\textbf{Falsification criteria.} We reject \(H_1\) if (i) mean accuracy
of B falls more than 15\% below A on the aggregate set, or (ii) blind
expert preference for B falls below 60\%.

\section{The AI Context OS
Architecture}\label{the-ai-context-os-architecture}

The framework comprises three layers:

\begin{enumerate}
\def\labelenumi{\arabic{enumi}.}
\tightlist
\item
  \textbf{Context Cores:} Curated markdown documents per application
  domain, containing system-prompt guidelines, key entities (classes,
  interfaces, APIs), and links to critical source files.
\item
  \textbf{Context Router:} A lightweight classifier (keyword- or
  embedding-based) that scores queries against cores.
\item
  \textbf{Dynamic Prompt Builder:} An orchestrator that compiles base
  instructions, active core guidelines, user preferences, and tool
  definitions.
\end{enumerate}

Figure 1 summarizes the compilation pipeline: the router scores the
query against cores; matched cores feed the prompt builder together with
base system instructions; the assembled prompt is sent to the LLM.

\textbf{Figure 1.} Context compilation pipeline: User Query
\(\rightarrow\) Context Router \(\rightarrow\) \{Workspace Core, Comm
Core, \ldots\} \(\rightarrow\) Dynamic Prompt Builder \(\rightarrow\)
LLM.

\subsection{Taxonomy of Context Cores}\label{taxonomy-of-context-cores}

We classify cores into five domains:

\begin{itemize}
\tightlist
\item
  \textbf{Personal/default core:} User metadata, preferences,
  conversational history.
\item
  \textbf{Workspace core:} File system, document stores, local database
  schemas.
\item
  \textbf{Communication core:} Messaging and mail integrations.
\item
  \textbf{System control core:} Execution paths, scripting, input
  automation.
\item
  \textbf{Browsing core:} Browsers, web sessions, tabs, URL requests.
\end{itemize}

\subsection{Core Construction}\label{core-construction}

Cores follow the \textbf{minimal context principle:}

\begin{enumerate}
\def\labelenumi{\arabic{enumi}.}
\tightlist
\item
  \textbf{Domain boundary mapping:} Identify distinct responsibilities
  (e.g., billing versus database migration).
\item
  \textbf{Entity extraction:} Capture public interfaces, settings maps,
  and structural invariants.
\item
  \textbf{Prose compression:} Replace raw listings with concise rules
  (e.g., \emph{``SQLite schema version 3.4''}).
\item
  \textbf{Exclusions:} State explicit out-of-scope boundaries to reduce
  hallucination.
\end{enumerate}

\subsection{Curation Cost}\label{curation-cost}

Context Cores introduce manual curation overhead. Across test projects:

{\def\LTcaptype{none} % do not increment counter
\begin{longtable}[]{@{}lcc@{}}
\toprule\noalign{}
Project & Cores & Authoring time \\
\midrule\noalign{}
\endhead
\bottomrule\noalign{}
\endlastfoot
MailAgent & 6 & \textasciitilde12 h \\
Django REST Framework & 5 & \textasciitilde18 h \\
Navorina & 11 & \textasciitilde40 h \\
Oiloop & 5 & \textasciitilde10 h \\
\end{longtable}
}

Initial overhead ranges from approximately 10 to 40 hours depending on
scale and coupling.

\section{Experimental Methodology}\label{experimental-methodology}

\subsection{Testbed}\label{testbed}

{\def\LTcaptype{none} % do not increment counter
\begin{longtable}[]{@{}lcccc@{}}
\toprule\noalign{}
Repository & LOC & Language & Questions & Cores \\
\midrule\noalign{}
\endhead
\bottomrule\noalign{}
\endlastfoot
MailAgent & 320k & TypeScript & 45 & 6 \\
Django REST Framework & \textasciitilde1.2M & Python & 42 & 5 \\
Navorina & \textasciitilde270k & Python & 42 & 11 \\
Oiloop & \textasciitilde180k & Swift (macOS) & 20 & 5 \\
\end{longtable}
}

\subsection{Design}\label{design}

We applied a within-subjects design with \texttt{gpt-4o-mini} as the
response and judge model:

\begin{itemize}
\tightlist
\item
  \textbf{Condition A:} All repository source and documentation.
\item
  \textbf{Condition B:} Routed core texts only.
\item
  \textbf{Condition C:} Static dependency subgraph (BFS expansion up to
  two hops).
\end{itemize}

\subsection{Metrics}\label{metrics}

\begin{itemize}
\tightlist
\item
  \textbf{Accuracy (0-3):} Correctness of technical details, frameworks,
  and APIs.
\item
  \textbf{Completeness (0-2):} Required steps, side effects, and file
  locations.
\item
  \textbf{Actionability (1-5):} Readability and executability without
  placeholders.
\item
  \textbf{Input tokens} and \textbf{latency (ms).}
\item
  \textbf{Hallucination rate (\%):} Incorrect APIs, classes, or files
  referenced.
\end{itemize}

\section{Results}\label{results}

Canonical runs: MailAgent \texttt{run-1781319187610}, Navorina
\texttt{run-1781143403051}, Oiloop Phase 3.1 \texttt{run-1781660908},
Django aggregate \texttt{run-drf-phase-2.1-aggregate} (per-question CSV
not preserved for Django).

\subsection{Cross-Project Comparison}\label{cross-project-comparison}

{\def\LTcaptype{none} % do not increment counter
\begin{longtable}[]{@{}
  >{\raggedright\arraybackslash}p{(\linewidth - 12\tabcolsep) * \real{0.1250}}
  >{\centering\arraybackslash}p{(\linewidth - 12\tabcolsep) * \real{0.0972}}
  >{\centering\arraybackslash}p{(\linewidth - 12\tabcolsep) * \real{0.1528}}
  >{\centering\arraybackslash}p{(\linewidth - 12\tabcolsep) * \real{0.1250}}
  >{\centering\arraybackslash}p{(\linewidth - 12\tabcolsep) * \real{0.1806}}
  >{\centering\arraybackslash}p{(\linewidth - 12\tabcolsep) * \real{0.1528}}
  >{\centering\arraybackslash}p{(\linewidth - 12\tabcolsep) * \real{0.1667}}@{}}
\toprule\noalign{}
\begin{minipage}[b]{\linewidth}\raggedright
Project
\end{minipage} & \begin{minipage}[b]{\linewidth}\centering
Cond.
\end{minipage} & \begin{minipage}[b]{\linewidth}\centering
Mean Acc.
\end{minipage} & \begin{minipage}[b]{\linewidth}\centering
Halluc.
\end{minipage} & \begin{minipage}[b]{\linewidth}\centering
Mean tokens
\end{minipage} & \begin{minipage}[b]{\linewidth}\centering
Router F1
\end{minipage} & \begin{minipage}[b]{\linewidth}\centering
CCR\_tokens
\end{minipage} \\
\midrule\noalign{}
\endhead
\bottomrule\noalign{}
\endlastfoot
MailAgent (45 Q) & A & 1.38 & 20\% & 88,113 & - & 1× \\
& B & \textbf{1.67} & 18\% & \textbf{10,615} & \textbf{1.00} &
\textbf{8.3×} \\
& C & 1.24 & 22\% & 20,955 & - & 4.2× \\
Django REST (42 Q)† & A & 1.35 & 18\% & \textasciitilde76,000 & - &
1× \\
& B & \textbf{1.68} & 22\% & \textbf{\textasciitilde2,000} & 0.72 &
\textbf{38×} \\
& C & 1.40 & \textbf{11\%} & \textasciitilde18,000 & - & 4.2× \\
Navorina (42 Q) & A & 1.00 & 24\% & 27,627 & - & 1× \\
& B & \textbf{1.19} & 19\% & \textbf{2,021} & 0.87 & \textbf{13.7×} \\
& C & 0.93 & \textbf{7\%} & 11,858 & - & 2.3× \\
Oiloop (20 Q, Ph 3.1) & A & 0.75 & 30\% & 76,663 & - & 1× \\
& B & \textbf{2.70} & \textbf{0\%} & \textbf{6,334} & 1.00 &
\textbf{12.1×} \\
& C & 2.35 & 10\% & 8,963 & - & 8.6× \\
\end{longtable}
}

† Django: aggregate export only; bootstrap CI not computed (no
\texttt{paired.csv}).

\subsection{Oiloop Phase 3.1}\label{oiloop-phase-3.1}

Phase 3.1 used \textbf{workspace-core v1.1.0}, corrected gold labels,
and gold routing (F1 = 1.0):

\begin{itemize}
\tightlist
\item
  B versus A accuracy: 2.70 versus 0.75 (Δ +1.95, N = 20).
\item
  Bootstrap 95\% CI on paired Δ (B − A): {[}+1.60, +2.30{]}; B exceeded
  A on 19/20 questions.
\item
  Production keyword router (Run 3): F1 = 1.0, B = 2.75, hallucination
  0\% (\texttt{run-prod-router-1781701118}).
\item
  Cross-cutting ablation (Run 2, 8 Q): B 2.875 \textgreater{} D (hybrid)
  2.50 \textgreater{} C 2.375 - \(H_{1f}\) rejected.
\end{itemize}

Prior Oiloop runs (pre-v1.1 cores) showed B 1.20 versus A 1.00; core
metadata richness (\(H_{1g}\)) explains the Phase 3.1 improvement.

\subsection{Compression and Latency}\label{compression-and-latency}

CCR\_tokens (mean tokens A / mean tokens B): 8.3× MailAgent, 38× Django,
13.7× Navorina, 12.1× Oiloop Phase 3.1. Condition B latency was
consistently lower than A on measured runs.

\subsection{Preference Pilots}\label{preference-pilots}

\textbf{Author-rated blind pilot} (10 questions, Phase 3.1 answers,
shuffled A/B): B preferred 7, equal 3, A 0 (n=10; not independent
raters; source: \texttt{research/human-blind-pilot-oiloop.csv}).

\textbf{Masked decode preference} (prior Oiloop canonical): B preferred
or equal on 75\% of 20 questions.

\section{Discussion}\label{discussion}

\subsection{Core Richness and Multi-Core
Routing}\label{core-richness-and-multi-core-routing}

Early Oiloop runs showed thin B accuracy (1.05-1.20) when cores lacked
enforcement chains and cross-cutting routes returned a single core.
Phase 3.1 addressed this via expanded metadata (entities, invariants,
rule-enforcement chains) and multi-core keyword routing (up to two cores
on tied scores). B accuracy rose from 1.20 to 2.70. Production default
on Oiloop is Condition B (\texttt{.fast} routing), not graph C.

\subsection{When Graph Retrieval Wins}\label{when-graph-retrieval-wins}

On Django REST and Navorina, C achieved lower hallucination than B (11\%
versus 22\%; 7\% versus 19\%) despite mixed accuracy. When trust
dominates, C is the recommended default for those profiles.

On Oiloop Phase 3.1, B exceeded C on accuracy and matched or beat C on
hallucination. Graph C remains an opt-in path (\texttt{.deep}).

\subsection{Hybrid Core plus Graph}\label{hybrid-core-plus-graph}

Condition D (routed core plus 15k-token graph supplement) did not match
C accuracy at lower cost than B on eight cross-cutting questions.
Multi-core B is the simpler strategy when cores are metadata-rich.

\subsection{Empirical Decision Matrix}\label{empirical-decision-matrix}

{\def\LTcaptype{none} % do not increment counter
\begin{longtable}[]{@{}
  >{\raggedright\arraybackslash}p{(\linewidth - 4\tabcolsep) * \real{0.3214}}
  >{\raggedright\arraybackslash}p{(\linewidth - 4\tabcolsep) * \real{0.3214}}
  >{\raggedright\arraybackslash}p{(\linewidth - 4\tabcolsep) * \real{0.3571}}@{}}
\toprule\noalign{}
\begin{minipage}[b]{\linewidth}\raggedright
Profile
\end{minipage} & \begin{minipage}[b]{\linewidth}\raggedright
Default
\end{minipage} & \begin{minipage}[b]{\linewidth}\raggedright
Evidence
\end{minipage} \\
\midrule\noalign{}
\endhead
\bottomrule\noalign{}
\endlastfoot
Narrow domain, F1 \(\geq\) 0.95 (MailAgent, Oiloop prod) & B & MailAgent
F1 = 1.0; Oiloop Run 3 F1 = 1.0 \\
Integrated OSS, hallucination-sensitive (DRF, Navorina) & C & C
hallucination \textless{} B \\
Integrated private app, rich cores (Oiloop Ph 3.1) & B & B 2.70;
multi-core wins cross-cutting \\
Full-repository dump & None & Dominated on cost and accuracy \\
\end{longtable}
}

\section{Limitations}\label{limitations}

\begin{enumerate}
\def\labelenumi{\arabic{enumi}.}
\tightlist
\item
  \textbf{LLM-as-judge:} All accuracy and hallucination scores use
  gpt-4o-mini; no inter-rater agreement with humans on the full
  149-question bank.
\item
  \textbf{Expert preference:} Masked decode and the author blind pilot
  are not substitutes for independent blind human studies.
\item
  \textbf{Django REST:} Only aggregate metrics exported; per-question
  \texttt{paired.csv} was not preserved.
\item
  \textbf{Curation cost:} 10-40 hours per codebase; semantic drift
  requires maintenance.
\item
  \textbf{Router dependency:} B quality scales with routing F1;
  mitigated by multi-core and semantic routing on DRF/Navorina.
\item
  \textbf{Single model:} All runs use gpt-4o-mini; generalization
  untested.
\item
  \textbf{Oiloop A baseline variance:} Condition A scores drifted
  between runs; Phase 3.1 is the reference.
\end{enumerate}

\section{Future Work}\label{future-work}

\begin{enumerate}
\def\labelenumi{\arabic{enumi}.}
\tightlist
\item
  Semi-automated core generation from AST and documentation.
\item
  Context drift monitors (PR hooks when source links change).
\item
  Full Django REST re-run with per-question export.
\item
  Independent human blind study (10-question pilot protocol published).
\item
  Multi-model replication on the Oiloop product path.
\end{enumerate}

\section{Conclusion}\label{conclusion}

Full-repository loading is token-inefficient for decision-scoped
questions. Domain-oriented Context Cores with multi-core routing deliver
8-38× compression (CCR\_tokens) on open-source projects and
\textasciitilde12× on Oiloop Phase 3.1. When cores are metadata-rich,
accuracy substantially exceeds monolithic dumps (Oiloop: B 2.70 versus A
0.75; bootstrap CI entirely above zero). Production defaults are
stratum-dependent: B on narrow, high-F1 paths; C on
hallucination-sensitive integrated OSS. Hybrid core-plus-graph did not
outperform rich multi-core B on Oiloop cross-cutting questions. We
release protocols, three full raw run exports, and an aggregate Django
export with stated limitations.

\section{References}\label{references}

{[}1{]} N. Liu et al., ``Lost in the Middle: How Language Models Use
Long Contexts,'' \emph{Transactions of the Association for Computational
Linguistics}, vol.~12, pp.~157-173, 2024.

{[}2{]} O. Press et al., ``Train Short, Test Long: Attention with Linear
Biases Enables Input Length Extrapolation,'' \emph{ICLR}, 2022.

{[}3{]} E. Evans, \emph{Domain-Driven Design: Tackling Complexity in the
Heart of Software}. Addison-Wesley, 2003.

{[}4{]} P. Lewis et al., ``Retrieval-Augmented Generation for
Knowledge-Intensive NLP Tasks,'' \emph{Advances in Neural Information
Processing Systems}, 2020.

{[}5{]} D. Edge et al., ``From Local to Global: A GraphRAG Approach to
Query-Focused Summarization,'' \emph{arXiv:2404.16130}, 2024.

{[}6{]} D. Guo et al., ``GraphCodeBERT: Pre-training Code
Representations with Data Flow,'' \emph{International Conference on
Learning Representations}, 2021.

{[}7{]} H. Jiang et al., ``LLMLingua: Compressing Prompts for
Accelerated Inference of Large Language Models,'' \emph{EMNLP}, 2023.

\section{Appendix A: Example Evaluation
Questions}\label{appendix-a-example-evaluation-questions}

\begin{itemize}
\tightlist
\item
  \textbf{OL01:} How does Oiloop store user preferences and chat
  history?
\item
  \textbf{OL05:} How does Oiloop handle sandboxed folder access on
  macOS?
\item
  \textbf{OL15:} How does EventKit Reminders authorization avoid
  freezing Oiloop in background mode?
\item
  \textbf{OL18:} How are Safari tab titles and URLs fetched?
\end{itemize}

\section{Appendix B: Context Core
Template}\label{appendix-b-context-core-template}

\begin{verbatim}
# [Core Name] Core

## Purpose
Decision domain description.

## Key Entities
- `EntityName`: Class/struct and responsibility.

## Invariants & Rules
1. Mandatory checks and constraints.

## Exclusions
- Explicit out-of-scope boundaries.
\end{verbatim}

\section{Appendix C: Router Sketch}\label{appendix-c-router-sketch}

\begin{verbatim}
struct ContextRouter {
    static func route(userInput: String) -> Set<ContextCore> {
        let lower = userInput.lowercased()
        // Score cores by keyword presence; return matching set
    }
}
\end{verbatim}

\section{Appendix D: Blinded Rater
Instructions}\label{appendix-d-blinded-rater-instructions}

\begin{enumerate}
\def\labelenumi{\arabic{enumi}.}
\tightlist
\item
  Answers are randomized as ``Option 1'' and ``Option 2.''
\item
  Rate each option on accuracy (0-3), completeness (0-2), and
  actionability (1-5).
\item
  Select the preferred option without knowing the retrieval source.
\end{enumerate}

\section{Data Availability}\label{data-availability}

Research repository:
\texttt{https://github.com/navorina-labs/AI-Context-OS}

Evaluation dataset: \texttt{research/questions.md}

Raw metrics: \texttt{context-os/evaluations/} and
\texttt{experiments/*/runs/}

\end{document}
